\makeabstract
{
Impact of Carbohydrate Sources On Prophage Induction In Phocaeicola vulgatus
}
{
Theint Theint Thu(1,2), Saria McKeithen-Mead(1), Jiawei Sun(1), Kerwyn C. Huang(1)
}
{
\textsuperscript{1}Department of Bioengineering, Stanford University, Stanford, CA 94305\\
\textsuperscript{2}Amherst College, Amherst, MA 01002
}
{
In a previous murine model of a stable synthetic gut community, a mobile genetic element (MGE) identified as a prophage transferred from the invading microbe P. vulgatus to a dominant community member. This transfer coincided with diet-dependent prophage induction, with increased copy number variation (CNV) observed in mice fed a standard diet but not a high-sugar diet. Here, we investigated how carbon source and concentration regulate prophage induction and lysis dynamics of P. vulgatus in vitro.
}
{
Follow-up in vitro experiments tested P. vulgatus growth in defined minimal media with various carbon sources such as D-glucose, L-arabinose, or rhamnogalacturonan-I (RG1) at concentrations ranging from 0.25–5 mg/mL.
}
{
Growth curves showed that at 5 mg/mL, lysis did not occur despite entry into the stationary phase, suggesting alternative regulatory pathways limit activation under saturating carbon conditions. Lysis was most pronounced at 2.5 mg/mL, and induction occurred even at 0.25 mg/mL. Phase contrast and LIVE/DEAD microscopy confirmed reduced viability post-lysis. qPCR targeting phage excision (attP) revealed prophage induction in a small subset of cells, with CNV increases preceding lysis.

}
{
These findings support that prophage induction in P. vulgatus is carbon source and concentration-dependent, and regulated by mechanisms beyond simple growth phase entry. Future work will involve expanding qPCR to measure phage particles for better understanding of what goes on at each time point in the growth assay, curing P. vulgatus of the prophage to assess fitness effects, and performing colony-forming unit assays to complement $OD_{600}$ measurements.
}

\newpage
